\documentclass[11pt,a4paper]{article}
\usepackage[utf8]{inputenc}
\usepackage[margin=1in]{geometry}
\usepackage{hyperref}
\usepackage{graphicx}
\usepackage{amsmath}
\usepackage{booktabs}

\title{RSCT Review: An interpretable prototype parts-based neural network for medical tabular data}
\author{Swarm-It Research Discovery\\\small Automated RSCT-Based Analysis}
\date{March 06, 2026}

\begin{document}
\maketitle

\begin{abstract}
This document provides an automated review of the paper ``An interpretable prototype parts-based neural network for medical tabular data'' from an RSCT (Representation-Solver Compatibility Testing) perspective. The paper achieved an RSCT relevance score of 6.2% and topic similarity of 40.1%.
\end{abstract}

\section{Paper Information}
\begin{itemize}
\item \textbf{Title:} An interpretable prototype parts-based neural network for medical tabular data
\item \textbf{Authors:} Jacek Karolczak, Jerzy Stefanowski
\item \textbf{Source:} \url{https://arxiv.org/abs/2603.05423v1}
\item \textbf{RSCT Relevance:} 6.2%
\item \textbf{Topic Match:} 40.1%
\item \textbf{Key Concepts:} alignment
\end{itemize}

\section{Original Abstract}
The ability to interpret machine learning model decisions is critical in such domains as healthcare, where trust in model predictions is as important as their accuracy. Inspired by the development of prototype parts-based deep neural networks in computer vision, we propose a new model for tabular data, specifically tailored to medical records, that requires discretization of diagnostic result norms. Unlike the original vision models that rely on the spatial structure, our method employs trainable patching over features describing a patient, to learn meaningful prototypical parts from structured data. These parts are represented as binary or discretized feature subsets. This allows the model to express prototypes in human-readable terms, enabling alignment with clinical language and case-based reasoning. Our proposed neural network is inherently interpretable and offers interpretable concept-based predictions by comparing the patient's description to learned prototypes in the latent space of the network. In experiments, we demonstrate that the model achieves classification performance competitive to widely used baseline models on medical benchmark datasets, while also offering transparency, bridging the gap between predictive performance and interpretability in clinical decision support.

\section{RSCT Analysis}
This paper presents research with 6% relevance to RSCT theory.

\textbf{Abstract Summary:} The ability to interpret machine learning model decisions is critical in such domains as healthcare, where trust in model predictions is as important as their accuracy. Inspired by the development of prototype parts-based deep neural networks in computer vision, we propose a new model for tabular data, specifically tailored to medical records, that requires discretization of diagnostic result norms. Unlike the original vision models that rely on the spatial structure, our method employs trainabl...

\textbf{RSCT Relevance:} The work touches on concepts related to alignment, which have connections to the Representation-Solver Compatibility Testing framework.

\textbf{Further Analysis:} A detailed manual review is recommended to fully assess the implications for RSCT-based certification approaches.


\subsection*{RSCT Certification Metrics}
\begin{tabular}{ll}
\textbf{Relevance (R):} & 0.375 \\
\textbf{Spurious (S):} & 0.319 \\
\textbf{Noise (N):} & 0.306 \\
\textbf{Kappa ($\kappa$):} & 0.495 \\
\end{tabular}

The kappa score of 0.495 indicates limited representation-solver compatibility.


\section{Relevance to Swarm-It}
This paper was identified by the Swarm-It research discovery pipeline as potentially relevant to RSCT-based AI certification. The combined score of 19.8% places it in the exploratory category for further investigation.

\vspace{1em}
\hrule
\vspace{0.5em}
\small{Generated by Swarm-It Research Discovery | \url{https://swarms.network} | RSCT Certified}

\end{document}
